\chapter{METODOLOGIA}

Neste estudo, foram construídas quatro aplicações distintas utilizando as linguagens de programação Java e Go, cada uma implementando dois protocolos de comunicação diferentes: HTTP e gRPC.
Essas aplicações foram projetadas com simplicidade e foco na comparação de desempenho, contendo apenas um \textit{endpoint}.
Esse \textit{endpoint} é responsável por receber uma lista de números aleatórios, iterar sobre essa lista e, em seguida, retornar o valor -1.
A escolha de uma operação simples permite isolar o desempenho da linguagem de programação e do protocolo de comunicação das complexidades inerentes à lógica de negócios, facilitando uma comparação direta e objetiva entre as tecnologias em teste.

Para cada combinação de linguagem de programação e protocolo, uma aplicação foi desenvolvida seguindo as melhores práticas e diretrizes específicas para cada linguagem e protocolo.

As aplicações foram executadas em um \textit{cluster} de kubernetes. (TODO: justificar porque)

Os testes foram executados em uma instalação do Ubuntu com a configuração mínima necessária, visando eliminar qualquer possível interferência de serviços ou aplicações secundárias no desempenho das aplicações em estudo.
Esse ambiente foi escolhido para garantir que os recursos do sistema fossem dedicados quase que exclusivamente às aplicações e aos processos de teste, proporcionando assim resultados mais precisos e confiáveis.

A métrica escolhida para avaliar o desempenho das aplicações neste estudo é o tempo de execução.
Essa métrica é fundamental para compreender a eficiência de cada combinação de linguagem de programação e protocolo de comunicação, permitindo mensurar o tempo em que cada aplicação processa e responde às requisições enviadas.
O tempo de execução é um indicador direto do desempenho e pode ser influenciado por diversos fatores, então será estudado o impacto das linguagens de programação e dos protocolos de comunicação nessa variável.

Para que seja possível entender ainda mais sobre os fatores que influenciam o estudo, dois experimentos serão executados com tamanhos de requisições denominados "\textit{small}" e "\textit{big}".
O \textit{payload} "\textit{small}", que compreende um arranjo de 200 números inteiros com valores aleatórios entre 0 e 1.000, foi meticulosamente calculado considerando a representação desses números em formato de texto.
O cálculo revelou que o tamanho médio de cada número é aproximadamente 2,91 bytes, e ao incluir uma vírgula como separador entre os números (exceto após o último número), o tamanho total estimado para o \textit{payload} "\textit{small}" é de aproximadamente 0.763 KB.
Este tamanho é representativo do tamanho médio de requisições HTTP comumente observadas em aplicações web, servindo como um \textit{benchmark} relevante para análises de desempenho \cite{chromium2012spdy}.

Por outro lado, o \textit{payload} "\textit{big}" \space é concebido para ser exatamente 1.024 vezes maior que o "\textit{small}", refletindo uma carga significativamente mais pesada. Isso eleva o tamanho do \textit{payload} "\textit{big}" para aproximadamente 
0,763
×
1.024
=
781,312KB. 
Essa escala ampliada é crucial para testar o desempenho da aplicação sob condições extremas de carga, permitindo uma investigação detalhada sobre como diferentes tamanhos de \textit{payload} afetam a latência, o \textit{throughput} e outros indicadores vitais de desempenho nas aplicações distribuídas. Tal abordagem garante uma compreensão abrangente do comportamento da aplicação em variados cenários de carga, fornecendo \textit{insights} valiosos para otimização e escalabilidade.

Observabilidade é a capacidade de inferir o estado interno de um sistema a partir de seus \textit{outputs}.
Esta definição é fundamental para o monitoramento e gerenciamento de sistemas complexos, como aplicações distribuídas em ambientes de Kubernetes \cite{Kang2022Observability}.
O Prometheus, por sua vez, é uma ferramenta \textit{open-source} amplamente reconhecida por sua eficácia na implementação de observabilidade, oferecendo funcionalidades como coleta de métricas em tempo real, processamento de consultas, visualização de dados e alertas~\cite{PrometheusDocumentation}. 
Dentro do ambiente Kubernetes, a kube-prometheus-stack é particularmente valorizada por integrar o Prometheus com o Grafana e outros componentes, facilitando assim a observabilidade ao prover um pacote completo que inclui coleta de métricas, visualização através de \textit{dashboards} e suporte adaptado ao Kubernetes.

Para a execução do experimento, foi desenvolvida uma aplicação de teste em Go, cujo propósito é automatizar o processo de teste, coleta e persistência de dados de desempenho das aplicações alvo.
O código principal inicia configurando o acesso ao \textit{cluster} Kubernetes. 

Para configurar e preparar o ambiente de teste, utilizamos o \texttt{helmfile}, uma ferramenta que facilita a orquestração de aplicações no Kubernetes por meio de arquivos de configuração YAML \cite{helmfileDocs}.

Após preparar o ambiente de teste com o helmfile, a aplicação espera os \textit{pods} estarem prontos, assegurando que todos os componentes necessários estejam operacionais antes de iniciar as requisições de teste.

A aplicação executa então uma série de requisições para as aplicações Java e Go, utilizando os protocolos HTTP e gRPC, de acordo com os parâmetros de experimento definidos. 

Após a conclusão das requisições, a aplicação coleta métricas de desempenho de cada uma das aplicações testadas, diretamente dos \textit{endpoints} do Prometheus configurados para monitoramento.
Finalmente, essas métricas são persistidas em um banco de dados SQL, permitindo análise posterior.
Esse processo é repetido para vários experimentos, garantindo uma ampla base de dados para análise de desempenho. 

Antes de dar início à execução prática do experimento, é imperativo estabelecer um projeto estatístico rigoroso e bem fundamentado.
Esse planejamento envolve a definição clara dos objetivos do estudo, a escolha do desenho experimental adequado e a determinação das técnicas estatísticas a serem aplicadas na análise dos dados.

Para as análises, as amostras não são pareadas devido à natureza intrínseca dos dados coletados, que envolvem comparações entre diferentes configurações de sistemas distribuídos sem correspondência direta entre as observações de cada grupo \cite{Illowsky2017IntroductoryStatistics}.
Esse método é adequado para situações onde os dados de um par não têm contraparte no outro, como é o caso ao comparar o desempenho de linguagens de programação e protocolos em ambientes distintos. 

Após a definição do tipo de amostra, torna-se crucial estabelecer um tamanho de amostra adequado que possa garantir a validade estatística dos resultados.
Para tanto, utiliza-se uma fórmula que permite definir o tamanho da amostra com base em um conjunto de observações preliminares.
Considerando um nível de confiança de 99\% e 9 graus de liberdade (10 observações preliminares), a variável~\(z\) tem o valor de 3,250 e representa o valor crítico da distribuição t de Student para esse nível de confiança e graus de liberdade especificados. O parâmetro \(r\) é definido como o erro tolerável, fixado em 5\%. A fórmula geral para o cálculo do tamanho da amostra (\(sz\)) para cada combinação de linguagem de programação e protocolo é dada por \cite{jain1991art}:
\[
sz = \left( \frac{100 \cdot z \cdot s}{r \cdot \bar{x}} \right)^2

\]
onde:
\begin{itemize}
    \item \(s\) é o desvio padrão das medições prévias,
    \item \(\bar{x}\) é a média das medições prévias,
    \item \(z\) é o valor da t de Student para o nível de confiança desejado,
    \item \(r\) é o erro percentual tolerável.
\end{itemize}

A tabela abaixo mostra os cáculos para cada aplicação
\begin{table}[h]
\centering
\begin{tabular}{lccc}
\hline
\textbf{Configuração} & \textbf{Média (\(\bar{x}\))} & \textbf{Desvio Padrão (\(s\))} & \textbf{Tamanho da Amostra (\(sz\))} \\
\hline
\multicolumn{4}{c}{\textbf{\textit{Big} Request}} \\
\hline
Java HTTP & 0.4178722 & 0.02489572 & 17.7724 \\
Java gRPC & 0.3943912 & 0.02710212 & 19.95165 \\
Go HTTP & 0.5219623 & 0.03239392 & 11.39082 \\
Go gRPC & 0.4666094 & 0.02128836 & 14.25365 \\
\hline
\end{tabular}
\caption{Determinando tamanho da amostra para requisições \textit{big}}
\label{tab:performance_big}
\end{table}

\begin{table}[h]
\centering
\begin{tabular}{lccc}
\hline
\textbf{Configuração} & \textbf{Média (\(\bar{x}\))} & \textbf{Desvio Padrão (\(s\))} & \textbf{Tamanho da Amostra (\(sz\))} \\
\hline
\multicolumn{4}{c}{\textbf{\textit{Small} Request}} \\
\hline
Java HTTP & 0.001589383 & 0.0002316207 & 190.3019 \\
Java gRPC & 0.001631577 & 0.0003373159 & 180.5865 \\
Go HTTP & 0.001399864 & 0.0004386 & 245.3175 \\
Go gRPC & 0.001186249 & 0.0005110324 & 341.624 \\
\hline
\end{tabular}
\caption{Determinando tamanho da amostra para requisições \textit{small}}
\label{tab:performance_small}
\end{table}

Conforme observado nas tabelas acima, o valor mais alto de tamanho de amostra necessário foi identificado para a configuração Go gRPC, exigindo 341 observações para alcançar o nível de confiança estatística necessário.
Esse requisito elevado pode ser atribuído ao desvio padrão relativamente maior observado nessa configuração, comparado às demais.
As outras configurações, devido a um desvio padrão menor, demonstraram a necessidade de um número menor de observações para atingir a mesma confiança estatística.
Seguindo uma abordagem estatística conservadora e com o objetivo de garantir a robustez dos resultados em todas as configurações, será definido o número de observações em 350 para todas as configurações.
Essa decisão assegura uma margem de confiança adequada, mesmo nas situações de maior variabilidade, e facilita a implementação e comparação dos resultados entre as diferentes tecnologias analisadas.

O experimento foi cuidadosamente executado e os dados relevantes foram coletados com precisão.
A primeira etapa na análise desses dados consistiu em verificar a aderência à normalidade da distribuição, um aspecto crucial para a determinação dos métodos estatísticos a serem aplicados posteriormente.
A normalidade dos dados é uma premissa fundamental para uma ampla gama de testes estatísticos.
Para avaliar a normalidade, serão empregadas duas abordagens complementares: a visualização dos dados por meio de histogramas e a aplicação do teste de Shapiro-Wilk.
O histograma oferece uma representação gráfica que pode sugerir visualmente a simetria e a forma da distribuição dos dados, enquanto o teste de Shapiro-Wilk fornece uma medida objetiva para avaliar a hipótese de normalidade.
Os resultados obtidos por meio do teste de Shapiro-Wilk são essenciais para esta análise, indicando de forma quantitativa se os dados se desviam significativamente de uma distribuição normal.
Abaixo, apresentamos uma tabela com os resultados do teste de Shapiro-Wilk para cada conjunto de dados coletado, permitindo uma avaliação precisa da normalidade e orientando as decisões sobre as técnicas estatísticas mais apropriadas para o prosseguimento da análise.

\begin{table}[h]
\centerin
\begin{tabular}{lccc}
\hline
\textbf{Configuração} & \textbf{\(W\)} & \textbf{Valor-p} \\
\hline
\multicolumn{3}{c}{\textbf{\textit{Small} Request}} \\
\hline
Java HTTP & 0.20075 & < 2.2e-16 \\
Java gRPC & 0.24176 & < 2.2e-16 \\
Go HTTP & 0.90576 & 5.994e-14 \\
Go gRPC & 0.93448 & 2.734e-11 \\
\hline
\end{tabular}
\caption{Resultados do Teste de Shapiro-Wilk para Requisições \textit{Small}}
\label{tab:shapiro-results-small}
\end{table}

\begin{table}[h]
\centering
\begin{tabular}{lccc}
\hline
\textbf{Configuração} & \textbf{\(W\)} & \textbf{Valor-p} \\
\hline
\multicolumn{3}{c}{\textbf{\textit{Big} Request}} \\
\hline
Java HTTP & 0.92997 & 9.419e-12 \\
Java gRPC & 0.95032 & 1.786e-09 \\
Go HTTP & 0.93732 & 5.497e-11 \\
Go gRPC & 0.92781 & 5.731e-12 \\
\hline
\end{tabular}
\caption{Resultados do Teste de Shapiro-Wilk para Requisições \textit{Big}}
\label{tab:shapiro-results-big}
\end{table}

A análise dos resultados obtidos pelos testes de Shapiro-Wilk revela que, independentemente do tamanho da requisição, nenhum dos conjuntos de dados segue uma distribuição normal. Os valores-p em todas as configurações são significativamente menores que o limiar de 0.05, o que leva à rejeição da hipótese de normalidade para todos eles. Mesmo nos casos onde os valores de \(W\) são relativamente mais altos, indicando uma possível proximidade com a normalidade, os valores-p sublinham uma clara descontinuidade com a distribuição normal.

Após a aplicação dos testes de Shapiro-Wilk, os quais revelaram um desvio da normalidade nos conjuntos de dados analisados, o próximo passo na análise consiste na avaliação visual dos histogramas.
Essa inspeção gráfica oferece uma perspectiva adicional sobre a distribuição dos dados, complementando os achados quantitativos dos testes estatísticos.
Os histogramas são especialmente valiosos para identificar características da distribuição, tais como assimetria e a presença de modas, que podem não ser evidentes através dos testes de normalidade.
Abaixo, encontram-se os histogramas para os dados associados às requisições de tamanho "\textit{small}".

% Histogramas para requisições Small
\begin{figure}[h]
    \centering
    \includegraphics[width=0.5\textwidth]{imgs/small_javahttp.png}
    \caption{Histograma Java HTTP com requisições \textit{small}}
    \label{fig:small_javahttp}
\end{figure}

\begin{figure}[h]
    \centering
    \includegraphics[width=0.5\textwidth]{imgs/small_javagrpc.png}
    \caption{Histograma Java gRPC com requisições \textit{small}}
    \label{fig:small_javagrpc}
\end{figure}

\begin{figure}[h]
    \centering
    \includegraphics[width=0.5\textwidth]{imgs/small_gohttp.png}
    \caption{Histograma Go HTTP com requisições \textit{small}}
    \label{fig:small_gohttp}
\end{figure}

\begin{figure}[h]
    \centering
    \includegraphics[width=0.5\textwidth]{imgs/small_gogrpc.png}
    \caption{Histograma Go gRPC com requisições \textit{small}}
    \label{fig:small_gogrpc}
\end{figure}

Continuando a análise, após a avaliação visual dos histogramas para as requisições de tamanho "\textit{small}", segue-se agora a inspeção dos histogramas correspondentes às requisições de tamanho "\textit{big}". 

% Histogramas para requisições Big
\begin{figure}[h]
    \centering
    \includegraphics[width=0.5\textwidth]{imgs/big_javahttp.png}
    \caption{Histograma Java HTTP com requisições \textit{big}}
    \label{fig:big_javahttp}
\end{figure}

\begin{figure}[h]
    \centering
    \includegraphics[width=0.5\textwidth]{imgs/big_javagrpc.png}
    \caption{Histograma Java gRPC com requisições \textit{big}}
    \label{fig:big_javagrpc}
\end{figure}

\begin{figure}[h]
    \centering
    \includegraphics[width=0.5\textwidth]{imgs/big_gohttp.png}
    \caption{Histograma Go HTTP com requisições \textit{big}}
    \label{fig:big_gohttp}
\end{figure}

\begin{figure}[h]
    \centering
    \includegraphics[width=0.5\textwidth]{imgs/big_gogrpc.png}
    \caption{Histograma Go gRPC com requisições \textit{big}}
    \label{fig:big_gogrpc}
\end{figure}

Após a análise dos histogramas, tanto para as requisições de tamanho "\textit{small}" quanto "\textit{big}", torna-se evidente que nenhuma das distribuições de dados segue uma normalidade perfeita. Os histogramas revelam desvios significativos em relação à forma de sino característica de uma distribuição normal, incluindo assimetrias e a presença de picos não característicos. Essa observação visual corrobora os resultados obtidos anteriormente pelos testes de Shapiro-Wilk, reforçando a conclusão de que as configurações analisadas apresentam distribuições que se afastam da normalidade

